\paragraph{\trDeterminantLine{}。}\label{sub:determinant-lines}
本节的讨论将主要围绕\trDeterminantLine{}展开，它是关联于\trFredholmOperator{}的一维实%
\trVectorSpace。实际上，令我们感兴趣的与其说是这样一条线本身，不如说只是它的两个可能%
\trOrientation{}所构成的集合。特别地，凡是说到两条\trDeterminantLine{}（或后文将要引入的
类似\trObject）之间的\trIsomorphism{}是\trCanonical{}的，其含义总是指：在乘以正实数的意义下
视为相同。做好了这一用语上的预防之后，设 $D$ 是两个实\trBanachSpace{}之间的一个%
\trFredholmOperator。$D$ 的\trDeterminantLine{}由\trKernel{}与\trCokernel{}的最高次\trExteriorPower{}作%
\trTensorProduct{}而成，更精确地说
\begin{equation*}
\det(D)=\lambda^{\mathrm{top}}\bigl(\mathrm{coker}(D)^{\vee}\bigr)\otimes
\lambda^{\mathrm{top}}\bigl(\ker(D)\bigr);
\end{equation*}
由于我们只关心\trOrientation，把\trCokernel{}取对偶本是多余的，但我们仍愿意保留这一写法，使得
这些公式与复\trDeterminantLine{}的公式保持可比。存在一个\trCanonical{}\trIsomorphism
\begin{equation}\label{eq:det-direct-sum}
\det(D_1)\otimes\det(D_2)\cong\det(D_1\oplus D_2),
\end{equation}
其定义方式是把诸向量都视为奇的，并使用 \trKoszul{}符号规则。显式地，若 $\{v_{k,i}\}$
是诸 $\ker(D_k)$ 的\trBasis，而 $\{w^{\vee}_{k,i}\}$ 是诸 $\mathrm{coker}(D_k)^{\vee}$ 的%
\trBasis，则该\trIsomorphism{}为
\begin{equation}\label{eq:det-koszul-iso}
\begin{split}
\bigl(\textstyle\bigwedge_i w^{\vee}_{1,i}\otimes\bigwedge_i v_{1,i}\bigr)\otimes
\bigl(\textstyle\bigwedge_i w^{\vee}_{2,i}\otimes\bigwedge_i v_{2,i}\bigr)
\qquad\qquad\qquad\quad\\
\longmapsto(-1)^{\dim(\mathrm{coker}\,D_2)\cdot\mathrm{index}(D_1)}
\textstyle\bigwedge_i w^{\vee}_{2,i}\wedge\bigwedge_i w^{\vee}_{1,i}\otimes
\textstyle\bigwedge_i v_{1,i}\wedge\bigwedge_i v_{2,i}.
\end{split}
\end{equation}
这是严格结合的，并且在\trGraded{}意义下也是交换的，意思是：上述两个\trIsomorphism{}的%
\trComposition
$\det(D_1)\otimes\det(D_2)\cong\det(D_1\oplus D_2)=\det(D_2\oplus D_1)\cong
\det(D_2)\otimes\det(D_1)$
与天真地交换因子的做法相差一个符号 $(-1)^{\mathrm{index}(D_1)\mathrm{index}(D_2)}$。
暂时先考虑有限维\trVectorSpace{}之间的\trLinear{}映射 $D\colon H\to K$。有唯一的一种
方式选取等同关系 $t_D\colon\det(D)\to\det(0)=\lambda^{\mathrm{top}}(K^{\vee})\otimes
\lambda^{\mathrm{top}}(H)$，使之对 $D$ 与左右两侧\trIsomorphism{}的\trComposition{}保持%
\trCompatible，并且满足以下三条附加规则：若 $D=0$，则 $t_D$ 是恒等映射；若
$D=\id\colon H\to H$ 而其中 $H$ 是一维\trVectorSpace，则
$t_D\colon\RR=\det(\id)\to H^{\vee}\otimes H$ 把 $1$ 映到 $v^{\vee}\otimes v$，其中
$\langle v^{\vee},v\rangle>0$；最后，诸\trIsomorphism{}与\trDirectSum{}\trCompatible，
其意义是下图可交换（诸 $0_k$ 是与 $D_k$ 相同的\trVectorSpace{}之间的零映射）：
\begin{equation*}
\begin{tikzcd}[column sep=6.5em]
\det(D_1\oplus D_2)\ar[r,"t_{D_1\oplus D_2}"]\ar[d,"\cong"']
  & \det(0_1\oplus 0_2)\ar[d,"\cong"]\\
\det(D_1)\otimes\det(D_2)\ar[r,"t_{D_1}\otimes t_{D_2}"']
  & \det(0_1)\otimes\det(0_2)
\end{tikzcd}
\end{equation*}
另外，对任意 $D\colon H\to K$ 也可以按下面的办法直接定义 $t_D$。取 $H$ 的一个\trBasis{}
$e_1,\ldots,e_m$ 与 $K$ 的一个\trBasis{} $f_1,\ldots,f_n$，使得
$D(e_1)=f_1,\ldots,D(e_k)=f_k$、$D(e_{k+1})=0,\ldots,D(e_m)=0$，并取相应的对偶\trBasis{}
$e_i^{\vee},f_i^{\vee}$。那么
\begin{equation}\label{eq:det-explicit-td}
\begin{split}
t_D:\ \bigl(f_n^{\vee}\wedge\cdots\wedge f_{k+1}^{\vee}\bigr)\otimes
\bigl(e_{k+1}\wedge\cdots\wedge e_m\bigr)&\in\det(D)\\
\longmapsto\bigl(f_n^{\vee}\wedge\cdots\wedge f_1^{\vee}\bigr)\otimes
\bigl(e_1\wedge\cdots\wedge e_m\bigr)&\in\det(0).
\end{split}
\end{equation}
有了这些准备，让我们回到一般情形：令 $\scr{F}(H,K)$ 为实\trBanachSpace{} $H$、$K$ 之间
全体\trFredholmOperator{}所成的空间。\trDeterminantLine{}\trLineBundle{}
$\underline{\det}\to\scr{F}(H,K)$ 以诸 $\det(D)$ 为\trFiber，其\trTopology{}结构可由以下性质刻画。
第一，若 $H$、$K$ 有限维，则诸\trIsomorphism{} $t_D$ 应当给出 $\underline{\det}$ 的一个%
\trTrivialization。第二，上文定义的映射
$\det(D_1)\otimes\det(D_2)\cong\det(D_1\oplus D_2)$ 是\trDeterminantLine{}\trLineBundle{}
之间的\trIsomorphism，定义在 $\scr{F}(H_1,K_1)\times\scr{F}(H_2,K_2)$ 之上。最后，若
我们限制到\trSurjective{}算子构成的子空间上，则等式
$\det(D)=\lambda^{\mathrm{top}}(\ker(D))$ 给出\trDeterminantLine{}\trLineBundle{}与核%
\trVectorBundle{}的最高次\trExteriorPower{}之间的一个\trIsomorphism（这是对经由%
\trStabilization{}定义\trDeterminantLine{}\trLineBundle{}的惯常做法的一个小小的变体）。
不用固定的\trBanachSpace，而是取两个\trBanach{}\trVectorBundle{}以及其间的一个%
\trFredholm{}映射，也同样可以得到基空间上的一个\trDeterminantLine{}\trLineBundle。
\begin{remark}\label{thm:det-sign-convention}
这里引入的符号约定取自 [87]（另见 [77, Appendix D]），在 [36, Volume I, pp.~62 and 362]
中被称为“观点 II”。提一下另一种可能的约定或许不无教益：把\eqref{eq:det-koszul-iso}
的右边换成
$\textstyle\bigwedge_i w^{\vee}_{1,i}\wedge\bigwedge_i w^{\vee}_{2,i}\otimes
\textstyle\bigwedge_i v_{1,i}\wedge\bigwedge_i v_{2,i}$（去掉符号），
并颠倒\eqref{eq:det-explicit-td} 中诸 $f_i^{\vee}$ 的次序。这就给出另一个%
\trLineBundle{} $\underline{\det}'$，其\trFiber{}与 $\underline{\det}$ 相同；两个丛%
\trIsomorphic，但并非以平凡的方式如此（恒等 $\underline{\det}\to\underline{\det}'$
不是连续映射）。
\end{remark}
